\section{Discussion}
\label{sec:discussion}

The main value of the benchmark is not that it names a winner; it constrains the decision space. Read in that way, the results support three operational regions rather than one champion list. First, EKF, RA-EKF, and UKF define the protection-oriented model-based region. They are not the cheapest estimators, but they remain the most credible broad-coverage choice when the disturbance class is uncertain and transient recovery matters as much as steady-state error. Second, SOGI-PLL, SOGI-FLL, and PLL define the low-cost synchronization region. They remain compelling when embedded implementation, deterministic execution, or ramp tracking dominates the design brief, but their behavior deteriorates more sharply under large discontinuities and composite IBR stress. Third, window-based methods such as TFT, IPDFT, Prony, and ESPRIT define the metering-oriented region: they can offer attractive steady-state precision or spectral behavior, but they pay for that strength through observation delay and, in several families, weaker protection-facing transient behavior.

Three practical decision rules follow. First, when the disturbance family is known in advance, the family-level table is the correct entry point because it preserves the relevant operating regime. Second, when the disturbance family is uncertain but the compute budget is moderate, the model-based filters are the safest broad-coverage choice in this study. Third, when the implementation budget is severe, low-cost loop methods remain credible only if the application can tolerate their weaker behavior under large discontinuities and composite IBR stress. These rules are more useful than a single leaderboard because they map the benchmark back to concrete protection, synchronization, and metering choices.

The statistical comparison layer strengthens that interpretation, but it should be read with care. The omnibus tests do not certify individual winners, and they do not replace the family-level tables. Their role is narrower and still important: they reject the idea that estimator-family labels are incidental metadata once the scenario summaries are pooled. That is enough to justify a reporting frame that honors family structure instead of flattening it.

One subtle implication concerns the role of headline winners. PLL leads the composite IBR multi-event family on RMSE, yet it is not the most robust global method. Koopman leads the isolated frequency-step case, yet that result does not overturn the broader pattern. ZCD enters the Pareto set, yet it does so for a reason that is mathematically true and operationally weak: vanishing CPU time paired with very poor error. A careful journal benchmark must distinguish these statements. Winning a local subproblem, surviving a Pareto screen, and offering a robust deployment recommendation are not equivalent achievements.

The older standards-style families also need to be interpreted in context. Clean ramps and magnitude steps remain necessary because they expose structural latency and tuning behavior under controlled conditions. However, they are not sufficient because they understate the failure modes that dominate low-inertia IBR operation: large phase discontinuities, noisy ringdown, and multi-event stress. The benchmark therefore supports a layered validation logic: standards-style cases establish baseline conformance, while IBR-specific and composite cases determine whether an estimator remains trustworthy when the signal assumptions behind that conformance are violated.
